% LAB BOREDOM — Section II appendices. GUIDED DRAFTING. Session 2026-08-21.
% Plan: plans/lab-boredom-appendices-plan-2026-08-21.md (all Phase-0 rulings received).
% ★ NUMBERING PROVISIONAL: these are SECTION-SCOPED designations (A.1/A.2/B/C/D, E optional);
%   the author will create other dissertation appendices later and the scheme WILL change
%   (memory: lab-boredom-appendix-renumbering). Every in-prose pointer stays ****** until
%   the dissertation-wide scheme exists — the designations below are working labels.
% Standing rulings: technical register permitted (D-G) · film-comparison tables EXCLUDED,
%   equipment-QC tables kept (D-B) · deployment-transfer topic DROPPED (D-C) · administered
%   survey wording is CANONICAL, no variant footnote (D-H) · papers bound as PDFs (D-E) ·
%   episode-only · data-as-agent · never "laboratory" · tree convention · no "ones".
% Tables: generated by analysis/appendix_tables.py -> out/appendix/D*.tex (spot-verified).
%
% BATCH A (headnote for A.1/A.2). REV 1 (2026-08-21, author note): "reanalyzed" ->
%   "reanalyzed and expanded upon."
% Verified facts: P1 authors King, Feeney, Tang, Das, Salvo (ASEE 2023 #37129); P2 authors
%   King, Lo, Das, Salvo (ASEE 2024 #44685, session-log-verified); IRB = UCI IRB Exempt
%   No. 20232678 (the 2024 form, per standing ruling); 3-of-8 pilot cohort fact; same
%   recordings reused in 2024; own publications in FIRST person.

Appendix A reproduces, in their published form, the two studies whose data this section
reanalyzed and expanded upon. Appendix A.1 is our 2023 pilot report, ``Work in Progress:
Physiological Assessment of Learning in a Virtual Reality Clinical Immersion Environment''
(King, Feeney, Tang, Das, and Salvo, ASEE 2023), which presented three of the first round's
eight participants. Appendix A.2 is our 2024 study, ``Assessment of Student Engagement in VR
Clinical Immersion Environments through Eye Tracking'' (King, Lo, Das, and Salvo, ASEE 2024),
which reported all twelve participants, drawing for the first round on the same recordings
the pilot had used. The analyses of this section pool those same twelve. Both studies were
conducted under UCI IRB Exempt No. 20232678.

% ============================================================================
% BATCH B (Appendix B — the survey instrument). Drafted 2026-08-21, awaiting author review.
% The seven questions are transcribed from the survey-response files and are IDENTICAL,
% word for word, across all twelve participants' records (cohort audit 2026-08-21, both
% rounds, one wording set). The seventh item's wording — open since L2 drafting — is hereby
% recovered and the ****** it carried resolves at integration.
% O-15 SLOT: the IRB-facing administration sentence is the author's, supplied at this
% batch's review — marked ****** below.
% Parsing counts verified against the tree 2026-08-21: exact 239 / converted 5 / flagged 1 /
% semantic 7 / missing 0 = 252. (Approved L2 prose still carries the pre-D-39 counts —
% flagged for a separate small revision batch with the author's L2 wording review.)

Appendix B reproduces the survey as it was administered. The seven questions below are
transcribed from the survey records themselves and are identical, word for word, across all
twelve participants' files. They were put to each participant aloud by the proctor and the
answers written down as given. The first question was asked just before each film began; the
remaining six were asked in the break that followed the film, their answers referring to the
film just watched.

\begin{enumerate}
\item How tired are you right now? (1-9; 1 completely awake, 9 completely exhausted)
\item On a scale from 1-9, how bored did you get watching this video? (1 not at all, 9
completely bored)
\item On a scale from 1-9, how interested and engaged in the video were you? (1 not engaged
at all, 9 intensely engaged)
\item Roughly how many minutes into the video until you started to get bored?
\item Did you fall asleep/fight to stay awake? (1 completely awake, 9 fell asleep)
\item Relatively speaking, roughly how long did the video feel to you?
\item On a scale from 1-9, how tired are you after watching the video (1 completely awake, 9
completely exhausted)
\end{enumerate}

Two quantities are derived rather than asked. Felt duration becomes a ratio against the
film's actual seventeen minutes, and depletion is the seventh answer minus the first. Because
the questions were read aloud and the answers recorded as free text, every answer passed
through a documented reading, and two coding rules govern the file. An answer of NA on any
1--9 item means the film produced none of the asked-after quality and is read as 1, the
scale's floor. A blank or NA on the minutes question means the participant never became
bored; those episodes are kept as their own category rather than assigned a number, since on
a duration item zero would mean bored from the first second. The reading of all 252 answers
(36 episodes by 7 items) breaks down as follows, and every interpreted value is stored
beside the exact words the participant gave:

\begin{tabular}{lcl}
\hline
reading & answers & meaning \\
\hline
exact & 239 & a plain number in range, no interpretation needed \\
converted & 5 & a unit or range converted (``25 sec'' to 0.42 minutes; ``10--12'' to 11) \\
flagged & 1 & an in-range answer carrying an anomaly, logged with its annotation \\
semantic & 7 & a meaningful non-number (never became bored; NA read as the scale floor) \\
missing & 0 & --- \\
\hline
\end{tabular}

% ============================================================================
% ★ APPENDIX B APPROVED by author 2026-08-21 (with the administration passage as revised
%   after the author's Q1-timing statement). O-15 RESOLVED. The four L2 revision edits
%   approved the same day.
%
% BATCH C1 (Appendix C — the criterion method in full). Drafted 2026-08-21, awaiting review.
% Base: PROSE-METHODS-CRITERION-VALIDITY-v2-2026-08-05.md §A ("do not rewrite" — numbers and
% logic stand). The file PREDATES the governing episode-only ruling (2026-08-07) and the
% D-39/D-40/D-41 re-runs, so the following DEVIATIONS are made and flagged to the author:
%   C1-a ¶1's opening two sentences (the films-apart question "answered by the preceding
%        tests") DELETED — those tests are excluded from the dissertation.
%   C1-b ¶5 closure figures refreshed to the current tree: +0.46/-0.42 -> +0.48/-0.44
%        (two-decimal style preserved; D-40 moved them).
%   C1-c ¶5's pupil sentence loses its film-separation clause -> "Pupil dilation does not
%        move with either report."
%   C1-d ¶6's closing clause recast episode-frame: "...and that no other recorded channel
%        does."
%   C1-e O-14 RESOLVED as ruled 2026-08-06: multiplicity handled in a closing FOOTNOTE, not
%        computed.
%   C1-f the v2 file's §C objections stay OUT (its own instruction: defense reserve, not
%        prose); §D notes are drafting apparatus, not appendix content.
% Everything else is v2 §A verbatim.

Appendix C states the criterion method in full. The instruments were adopted because the
participants' own reports were treated as the criterion, and each channel was scored by
whether it agreed with what the participant said afterward. Restated in the form the
reprocessing can test: does a physiological channel move with a person's own account of their
boredom, episode by episode?

Both sides of the comparison are converted to within-subject standard scores before anything
is correlated. Each participant's value for one film is expressed as its distance from that
participant's own three-film mean, in units of that participant's own standard deviation. The
step is not cosmetic. Absolute pupil diameter, habitual gaze range and rating style vary
enough between twelve people that an uncentered correlation would chiefly measure who has a
wide-ranging eye or a generous hand with a nine-point scale. After centering, each participant
contributes only the shape of their three films relative to themselves, which is the quantity
the question is about. The association is then computed across all thirty-six
participant-by-film cells.

The reported coefficient is Spearman's rank correlation. It correlates the order of the values
rather than the values themselves, which suits this material for two reasons. No theory
predicts that self-reported boredom rises in proportion to millimeters of pupil or degrees of
gaze deviation, so a statistic that assumes a straight line assumes something the design does
not license; and with twelve participants a single unusual episode can bend a linear
coefficient in a way it cannot bend a rank one. Pearson's product-moment correlation is
computed alongside and reported in full in the tables that follow, both because the two
coefficients should be compared rather than one quietly chosen, and because their agreement is
itself informative. Across every channel and both criteria the median absolute difference
between them is 0.032 and the largest is 0.117. Where a monotonic relationship is also close
to linear, the rank statistic loses almost nothing, so the choice of Spearman costs little.
Nothing in it flatters the result: on every association that carries the finding the Pearson
coefficient is the \textit{larger} of the two in magnitude, so the reported figure understates
the association rather than inflating it.

The probability attached to each coefficient is obtained by permutation rather than from the
closed-form distribution, and the reason is a property of the design that the standard formula
cannot accommodate. Thirty-six paired observations are not thirty-six independent
observations: they are twelve participants contributing three films each. Within-subject
centering, which the comparison requires for the reasons given above, makes the dependence
structural rather than incidental---constraining three values to a mean of zero and a standard
deviation of one means that knowing two of them very nearly fixes the third. Both the rank and
the linear coefficient carry analytic p-values that assume independent observations, so both
are anti-conservative here. The correction is to build the null distribution from the data's
own structure: the three film labels are shuffled within each participant, never across
participants, the coefficient is recomputed, and the procedure repeats twenty thousand times
with a fixed seed. The proportion of shuffles producing an association at least as strong as
the observed one is the reported probability. Analytic values are retained beside the
permutation values throughout, so the size of the correction is visible rather than absorbed.

That correction is consequential, and it is applied to both coefficients, which shows it to be
a feature of the design rather than an artifact of ranking. Where an association is weak the
two probabilities barely differ; where it is strong the analytic value runs three to eight
times too small, so the distortion is largest exactly where a reader would lean on it hardest.
Three associations that clear the conventional threshold on the analytic value do not clear it
on the permutation value, and they are reported as not clearing it. What survives is stated
without hedging. On the per-participant baseline, the one assigned to this question, gaze
deviation moves with self-reported boredom at rho = +0.59 and against self-reported engagement
at rho = $-$0.59; the per-file baseline stands next at +0.52 and $-$0.53 and points the same
way, so the result does not turn on which of the two is used. Eye closure moves with reported
boredom at rho = +0.48 and against reported engagement at $-$0.44. All of these survive the
permutation test. Pupil dilation does not move with either report.

Two limits bound what any of this can be taken to show. The cohort is twelve, so these
coefficients describe this group and support no estimate of a population value; they are
reported for their direction and their relative ordering, not as parameters. And the
comparison is episode-level throughout---each physiological figure is an aggregate over a
whole seventeen-minute film and each rating is given in the break immediately afterward---so
nothing here time-locks a moment of experience to a moment of signal. What the analysis
establishes is that two of the eye channels track a person's own retrospective account across
their three films, and that no other recorded channel does.\footnote{The criterion tests are
reported as a family of descriptive associations, not as a set of independent hypothesis
tests, and the argument uses the ordering of channels rather than any individual probability.
No multiplicity correction is therefore applied to the family, by choice rather than
oversight.}
